% ---------------------------------------------------------------------------
% 7 Results
% Every number here comes from evidence/measurements.yaml, telemetry.yaml, or scaling.yaml.
% The complete-model timings come from the dedicated-node run of 2026-08-12 (n = 1).
% Per-block timings sourced only to the shared-partition walkthrough remain withheld.
% ---------------------------------------------------------------------------

\section{Results}
\label{sec:results}

The client-assisted circuit defined in Section~\ref{sec:protocol} is evaluated along four axes:
numerical agreement, the executed operation schedule, memory, and the cost of a complete encrypted
inference. The complete model has been executed once on an uncontended node, and that single
execution is the source of every timing reported here. Per-block timings from earlier
shared-partition runs are not reported, because host contention cannot be separated from the
measurement after the fact.

\subsection{Correctness}

Correctness rests on the two-stage chain described in Section~\ref{sec:baseline}. Before encrypted
evaluation, the independent NumPy reference is asserted against the upstream PyTorch computation at
a tolerance of $2\times10^{-5}$, separately for all $12$ transformer blocks and the classifier
head; export aborts on a divergence. The encrypted circuit is then compared with that validated
reference. This ordering prevents an unchecked reimplementation from serving as the sole judge of
its encrypted counterpart.

For one block at the $103$-token task prompt, the decrypted output has relative infinity-norm error
$4.64\times10^{-9}$ against an acceptance tolerance of $4\times10^{-2}$ fixed before the run,
leaving seven orders of magnitude of headroom. The worst token-relative error is
$4.7\times10^{-9}$, and the largest absolute deviation is $7.44\times10^{-8}$. These values
establish numerical agreement for the encrypted block with the released weights and exact client
nonlinearities.

The complete model is held to the same unchanged tolerance. Evaluating all $12$ released blocks
and the released task head at the $103$-token task prompt yields a decrypted classification margin
that differs from the plaintext oracle margin by $8.56\times10^{-9}$ in relative terms and
$1.05\times10^{-7}$ in absolute terms, and the decrypted class label matches the frozen reference
label.

Error does accumulate with depth, and the size of that accumulation is the useful result. The
hidden-state relative infinity norm is $4.09\times10^{-9}$ after the first block and
$2.07\times10^{-8}$ after the twelfth, while the largest absolute deviation grows from
$6.56\times10^{-8}$ to $7.97\times10^{-6}$, roughly two orders of magnitude across the model.
Growth of that rate leaves the twelfth-block activation six orders of magnitude inside the
acceptance tolerance, and leaves the margin the classifier actually acts on inside it by nine. The
claim is therefore not that approximation error fails to accumulate, but that it accumulates far
too slowly over this depth to change the decision.

\subsection{Executed work}

\begin{table}[t]\centering\small
\caption{Checked operation schedule for the evaluated transformer block.}
\label{tab:cost}
\begin{tabularx}{\textwidth}{Y r l}
\toprule
Operation & Count & Executor \\
\midrule
Dense products & 156 & compute provider \\
Ciphertext--plaintext multiplications & 177{,}734 & compute provider \\
Ciphertext--ciphertext multiplications & 1{,}506 & compute provider \\
Rotations & 8{,}173 & compute provider \\
Slot reductions & 8{,}776 & compute provider \\
Physical client boundary crossings & 857 & data owner \\
Logical nonlinear values & 129{,}162 & data owner \\
\bottomrule
\end{tabularx}
\end{table}

The implementation asserts every count in Table~\ref{tab:cost}, the boundary sequence, and the
remaining multiplicative depth before the MLP. Any deviation aborts the run. This turns the
operation schedule into evidence that the optimized implementation executes the same circuit, not
an assurance inferred from a passing output.

Packing separates physical interaction from nonlinear model work. The $857$ physical client
boundary crossings carry $129{,}162$ logical values. Softmax accounts for $128{,}544$ of those
values, or $99.5\%$; LayerNorm and GELU together account for the remainder. Boundary optimization
is therefore principally an attention problem, even though the protocol contains three
nonlinearity classes.

The dense schedule also records the benefit of token packing. Eight-token packing reduces the
serial-equivalent dense-product count from $1{,}236$ to $156$, a $7.92\times$ reduction in
operations. Rotations, masking, client work, and host preparation remain in the packed circuit, so
the factor describes operation count rather than latency.

\subsection{Memory}

The evaluating process peaks at $9839$~MiB of GPU memory, within the $40$~GB device used for the
experiment. This is the evaluating process footprint rather than device-wide occupancy.

Host memory follows the lifetime of encoded plaintext templates. Query/key/value projection reaches
$25{,}530$~MiB of target-process resident memory. A cache flush precedes attention, whose recorded
resident memory is $12{,}434$~MiB. The MLP holds the largest live weight set and reaches the block
peak of $48{,}610$~MiB. These stage transitions match the explicit flush points in
Section~\ref{sec:optimization}; encoded weights do not accumulate across the entire block.

The two memory results have different mechanisms. The non-interactive composition configuration
exceeded its tested $80$~GB GPU envelope while loading the deep context and evaluation keys.
Client re-encryption shortens the encrypted segments and resets their depth budget, moving that
accelerator-memory barrier. Stage-bounded caching prevents the resulting host-side plaintext
templates from becoming an unbounded second barrier.

\begin{figure}[tb]\centering
  \includegraphics[width=\textwidth]{fig_timeline.pdf}
  \caption{Resource trace of the complete encrypted inference. GPU utilization never exceeds $51\%$, so the evaluation is not accelerator-bound even at whole-model scale. Host memory oscillates once per block, following the stage flushes of Section~\ref{sec:optimization}, rather than accumulating across the twelve blocks; GPU memory reaches its peak during the first block and does not grow with block index. The twelve cycles are the twelve blocks.}
  \label{fig:timeline}
\end{figure}

The complete model confirms that both bounds hold at full depth. GPU memory rises during context
and key loading, reaches $9839$~MiB after roughly $400$~s, and then holds that value exactly for
the remainder of the run: the same footprint the single block reaches, with no growth in block
index; host resident memory peaks at $50{,}930$~MiB, consistent with the same MLP-stage maximum
recurring once per block rather than accumulating. The memory barrier that stopped the
non-interactive configuration evaluated here is therefore not deferred by client assistance but
removed: twelve blocks cost what one block costs.

\subsection{Composition and evaluation boundary}

\begin{table}[t]\centering\small
\caption{Encrypted evaluation completed at each reported scope. Hidden-state error is a
relative infinity norm over the activation; margin error is the scalar the classifier acts on.
The two are different quantities and are kept in separate columns.}
\label{tab:scope}
\begin{tabularx}{\textwidth}{Y l l l}
\toprule
Scope & Status & Hidden-state error & Margin error \\
\midrule
One block, 2 tokens & measured & --- & --- \\
Two consecutive blocks, 2 tokens, client refresh & measured & $8.48\times10^{-11}$ & --- \\
One block, 103 tokens (task length) & measured & $4.64\times10^{-9}$ & --- \\
Twelve blocks and task head, 103 tokens & measured & $2.07\times10^{-8}$ & $8.56\times10^{-9}$ \\
Encrypted token-index lookup & out of scope & --- & --- \\
\bottomrule
\end{tabularx}
\end{table}

Two consecutive released blocks compose at reduced sequence length through a full-hidden-state
client refresh, reaching $8.48\times10^{-11}$ relative error after the second block. GPU memory
peaks in the first block and does not increase in the second. This experiment checks the mechanism
used for deeper composition: a fresh encryption starts the next block at level~$0$ without
carrying forward the preceding block's multiplicative depth or GPU-resident temporaries.

The complete model exercises that mechanism eleven more times. All $12$ blocks and the task head
run in one key lineage, separated by $11$ full-hidden-state refreshes between consecutive blocks
and one further refresh of the final token into the head. Multiplicative depth stays at $13$
throughout, the same budget the single block requires, and the run performs no homomorphic
bootstrap at any point. Depth is thus a property of the longest segment between two client
boundaries rather than of network depth, which is what makes a twelve-block encrypted lineage
expressible at all under these parameters.

\subsection{Cost of a complete encrypted inference}

\begin{figure}[tb]\centering
  \includegraphics[width=\textwidth]{fig_cost_split.pdf}
  \caption{How one complete encrypted inference divides between the two parties. The compute provider performs the encrypted linear algebra on one A100; the data owner evaluates the nonlinearities on CPU with no GPU. Server and client shares sum exactly to encrypted evaluation. Single execution on an uncontended node.}
  \label{fig:costsplit}
\end{figure}

\begin{table}[t]\centering\small
\caption{One complete encrypted inference: twelve blocks and the released task head at the
$103$-token task prompt, on one A100 GPU node using 32 CPU threads. Single execution on an
uncontended node; these are one sample's values, not averages over repetitions. Shares are rounded
independently and need not sum to $100\%$.}
\label{tab:fullcost}
\begin{tabularx}{\textwidth}{Y r r}
\toprule
Phase & Seconds & Share of wall \\
\midrule
Load released weights and verified reference & 23.494 & 0.4\% \\
Crypto context, key generation, GPU key upload & 3.945 & 0.1\% \\
Encrypt the embedded input & 1.265 & 0.0\% \\
\textbf{Encrypted evaluation} & \textbf{6594.748} & \textbf{98.7\%} \\
\quad server linear algebra (GPU) & 5638.561 & 84.4\% \\
\quad client boundaries (CPU, exact) & 956.187 & 14.3\% \\
Final decrypt & 0.044 & 0.0\% \\
Uninstrumented remainder & 59.504 & 0.9\% \\
\midrule
\textbf{Wall clock} & \textbf{6683} & \\
\bottomrule
\end{tabularx}
\end{table}

A complete encrypted inference takes $6683$~s of wall clock, or $1.86$~hours. Encrypted evaluation
accounts for $6594.748$~s of that, and setup is negligible at this scale: context construction and
key generation together take under four seconds, and the single final decryption takes $0.044$~s.
The instrumented phases are a decomposition rather than a partition, leaving $59.504$~s, or $0.9\%$
of wall clock, outside them; Table~\ref{tab:fullcost} shows that remainder rather than distributing
it silently.

Within encrypted evaluation the split between the two parties is exact: server linear algebra uses
$5638.561$~s and the data owner's boundary work uses $956.187$~s, and the two sum to the evaluation
total without residue. The compute provider therefore carries $85.5\%$ of the encrypted work and
the data owner $14.5\%$, with the data owner's share running on CPU alone. Across $10{,}299$
physical boundary crossings the mean crossing costs $0.093$~s, and those crossings carry
$1{,}551{,}081$ individual nonlinear values, so the interaction count understates the nonlinear work
by more than two orders of magnitude.

The twelve blocks account for $6565.993$~s and the encrypted task head for $9.046$~s, so the head
is not a meaningful part of the cost. The twelve refreshes together cost $19.710$~s, which places
the price of the composition mechanism at roughly $0.3\%$ of encrypted evaluation. Dividing block
evaluation by twelve gives $547.166$~s per block. The twelve individual block times span
$533.869$~s to $558.943$~s, a spread of $4.7\%$ across architecturally identical blocks on one
node. That bounds within-run scheduling jitter; it is not run-to-run variance, which a single
execution cannot report.

\subsection{Sequence-length structure}

\begin{figure}[tb]\centering
  \includegraphics[width=\textwidth]{fig_scaling.pdf}
  \caption{Circuit size across the encrypted-scope tasks. The bars show causal attention score tiles per block, which follow from the ciphertext-group count and not from any timing model; only genomic signal recognition has a measured complete encrypted inference. Token counts follow from each task's window under 6-mer tokenization.}
  \label{fig:scaling}
\end{figure}

Sequence length follows from each task format. Genomic-signal recognition uses $600$ base pairs,
which become $100$ sequence tokens under $6$-mer tokenization; three task-format tokens produce the
$103$-token encrypted anchor. Eight-token packing maps that prompt to $13$ ciphertext groups. The
shorter promoter and splice-site prompts use fewer groups under the same layout, yielding smaller
causal schedules.

The abundance-regression task lies outside the encrypted scope. Its $10{,}500$ base pairs produce
$1{,}755$ tokens and $220$ ciphertext groups under the same tokenization rule. The causal attention
schedule grows quadratically in the number of groups, placing this task in a different circuit-size
regime from the classification prompts.

\subsection{Feasibility}
\label{sec:verdict-feasibility}

Complete encrypted inference for DNAGPT is feasible under the stated protocol. All $12$ released
blocks and the released task head evaluate at the task's own prompt length in a single key lineage,
and the decrypted prediction reproduces the frozen plaintext label with $8.56\times10^{-9}$
relative margin error against a tolerance fixed before the run. The claim is no longer an argument
from a block plus a composition mechanism; the complete computation has been run and accepted.

Three properties make that outcome more than a single passing execution. Multiplicative depth stays
at $13$ across the whole model, so depth is bounded by the longest segment between client
boundaries and not by network depth. GPU memory holds at $9839$~MiB once loading completes, rather than
growing with block index, so composition does not accumulate accelerator state. And no
homomorphic bootstrap is performed anywhere, which is what removes the deep-context and
evaluation-key residency that stopped the non-interactive configuration evaluated in
Section~\ref{sec:protocol}. Correctness, depth, and accelerator memory are therefore all resolved
for a complete model at task length. Whether they would have obstructed a different
non-interactive implementation is not settled by this study, and Section~\ref{sec:related} reports
prior systems that do not encounter the same envelope.

What remains outside this verdict is the encrypted token-index lookup, which the protocol does not
attempt, and the transport layer, which is evaluated in one process rather than across a network.
Section~\ref{sec:limitations} states both.

\subsection{Practicality, judged separately}
\label{sec:verdict-practicality}

Feasibility does not make the system practical, and the measured cost says so plainly. One
encrypted classification of one $600$-base-pair sequence occupied an A100 node for $1.86$~hours,
against a plaintext forward pass of the same released model that completes in well under a second
on commodity hardware. That rules out interactive per-query use outright. It does not rule out batch or offline analysis of
a sample set where disclosing the sequence is unacceptable and the alternative is not analysing it
at all, since in that setting hours per sequence is a cost rather than a barrier.

The comparison that decides the verdict, however, is not against plaintext inference but against
other encrypted systems, and it does not favour this design. THOR reports $602$~s for a
non-interactive encrypted forward pass of a $12$-layer, $768$-wide encoder at $128$ tokens on a
single A100, at the ring degree, multiplicative depth, and security level used here, with no client
participation at any point~\cite{thor}; MOAI reports a further reduction against
that~\cite{moai2025}. The $6683$~s measured here is roughly an order of magnitude above it at a
shorter prompt. The practicality verdict is therefore negative in two distinct senses: too slow for
interactive use in absolute terms, and behind the non-interactive state of the art in relative
terms. Section~\ref{sec:related} sets out what the design offers in exchange, which is exactness,
freedom from bootstrapping, and an output verifiable against a plaintext reference, rather than
speed.

Two qualifications bound how far our own number should be carried, and both of them widen rather
than narrow the gap above. It is a single execution, so it supports a statement about the cost of
this inference and not a distribution, a mean, or a service rate; repeated runs would be needed
before quoting a latency target. And it is measured in one process, so it excludes the serialization
and network transport that a deployed service would add to every one of the $10{,}299$ boundary
crossings, whereas a fully non-interactive system incurs no such transport at all. The comparison is
therefore generous to the design evaluated here.

The measurement also indicates that this cost is not a floor imposed by cryptography. GPU
utilization holds a median of $32\%$ and a maximum of $51\%$ across the complete run, so the
evaluation is not limited by accelerator throughput even at whole-model scale; the cost sits in
host-side encoding and in the sequential client boundary, which
Section~\ref{sec:optimization} identifies as addressable. Feasibility and practicality thus resolve
differently and for different reasons: the encrypted path is correct and complete, while its
present latency confines it to batch and offline use.
